\documentclass[11pt,a4paper]{article}

% ============================================================
%  Metodología completa del último entrenamiento (split 361/63/84, 100 épocas):
%  ingesta de datos, arquitecturas capa por capa, pesos de partida, inflado 2D -> 3D,
%  resultados esperados y cómo se calcula cada número.
%  Figuras nuevas: python scripts/5_reports/build_report_metodologia.py
%  (las demás vienen de build_report_split361.py, report_split361_100.py y
%  analyze_latent_split361.py; ver el apéndice).
%  Compilar desde docs/: pdflatex informe_metodologia_split361.tex (dos veces)
% ============================================================

\usepackage[a4paper,margin=2.2cm]{geometry}
\usepackage[T1]{fontenc}
\usepackage{lmodern}
\usepackage[spanish,es-noshorthands]{babel}
\usepackage{microtype}
\usepackage{xcolor}
\usepackage{amsmath,amssymb}
\usepackage{booktabs}
\usepackage{array}
\usepackage{tabularx}
\usepackage{multirow}
\usepackage{graphicx}
\usepackage{caption}
\usepackage{float}
\usepackage{placeins}
\usepackage{tocloft}
\usepackage{enumitem}
\usepackage{fancyhdr}
\usepackage{titlesec}
\usepackage[most]{tcolorbox}
\usepackage{tikz}
\usetikzlibrary{arrows.meta,positioning,fit,backgrounds,calc}

\definecolor{primaryblue}{HTML}{1B4F72}
\definecolor{steelblue}{HTML}{2E86AB}
\definecolor{accentorange}{HTML}{E67E22}
\definecolor{softgray}{HTML}{85929E}
\definecolor{panelbg}{HTML}{EAF2F8}
\definecolor{warnbg}{HTML}{FDEBD0}
\definecolor{tldrbg}{HTML}{EBF5FB}
\definecolor{vicregc}{HTML}{2A78D6}
\definecolor{byolc}{HTML}{EB6834}
\definecolor{dinoc}{HTML}{1BAF7A}
\definecolor{supc}{HTML}{52514E}
\definecolor{okc}{HTML}{1E8449}
\definecolor{noc}{HTML}{B03A2E}
\definecolor{midc}{HTML}{B9770E}

\usepackage{xurl}
\usepackage[colorlinks=true,linkcolor=primaryblue,citecolor=primaryblue,urlcolor=steelblue]{hyperref}

\newtcolorbox{tldrbox}[1][En una frase]{colback=tldrbg,colframe=primaryblue,coltitle=white,
  colbacktitle=primaryblue,fonttitle=\bfseries,boxrule=0.9pt,arc=3pt,
  left=8pt,right=8pt,top=5pt,bottom=5pt,title=#1}
\newtcolorbox{ideabox}[1][Idea clave]{colback=panelbg,colframe=steelblue,coltitle=white,
  colbacktitle=steelblue,fonttitle=\bfseries,boxrule=0.8pt,arc=2pt,
  left=8pt,right=8pt,top=5pt,bottom=5pt,title=#1}
\newtcolorbox{warnbox}[1][Ojo con esto]{colback=warnbg,colframe=accentorange,coltitle=white,
  colbacktitle=accentorange,fonttitle=\bfseries,boxrule=0.8pt,arc=2pt,
  left=8pt,right=8pt,top=5pt,bottom=5pt,title=#1}

\pagestyle{fancy}
\fancyhf{}
\renewcommand{\headrulewidth}{0.5pt}
\fancyhead[L]{\small\color{softgray}Metodología · split 361/63/84 · 100 épocas}
\fancyhead[R]{\small\color{softgray}VICReg / BYOL / DINOv2}
\fancyfoot[C]{\small\thepage}
\setlength{\headheight}{14pt}

\titleformat{\section}{\normalfont\Large\bfseries\color{primaryblue}}{\thesection}{0.6em}{}
\titleformat{\subsection}{\normalfont\large\bfseries\color{steelblue}}{\thesubsection}{0.6em}{}
\titleformat{\subsubsection}{\normalfont\normalsize\bfseries\color{primaryblue}}{\thesubsubsection}{0.6em}{}

\graphicspath{{figures_split361/}}
\setlength{\parskip}{4pt}
\setlength{\parindent}{0pt}
\setlength{\emergencystretch}{3em}
\setlist[itemize]{topsep=2pt,itemsep=2pt}
\setlist[enumerate]{topsep=2pt,itemsep=2pt}
\captionsetup{font=small,labelfont={bf,color=primaryblue}}
\renewcommand{\arraystretch}{1.18}
\setlength{\cftsecnumwidth}{2.2em}
\setlength{\cftsubsecnumwidth}{3.0em}

\newcommand{\vox}[1]{$#1^3$}
\newcommand{\si}{\textcolor{okc}{\textbf{sí}}}
\newcommand{\no}{\textcolor{noc}{\textbf{no}}}
\newcommand{\parcial}{\textcolor{midc}{\textbf{en parte}}}
\newcommand{\sig}{$^{*}$}

\tikzset{
  stage/.style={rectangle, rounded corners=2pt, draw=primaryblue, thick, fill=panelbg,
    text width=2.55cm, align=center, minimum height=1.25cm, inner sep=3pt, font=\scriptsize},
  stagew/.style={stage, fill=white},
  stageo/.style={stage, draw=accentorange, fill=accentorange!8},
  arr/.style={-{Latex[length=2mm]}, thick, primaryblue},
  arrg/.style={-{Latex[length=2mm]}, thick, softgray},
  lane/.style={rectangle, rounded corners=2pt, draw=#1, thick, fill=#1!7, text width=3.35cm,
    align=left, minimum height=2.9cm, inner sep=4pt, font=\scriptsize, anchor=north},
}

\begin{document}

% ============================================================
\begin{titlepage}
\centering
\vspace*{0.8cm}
{\Huge\bfseries\color{primaryblue} SSL 3D en OASIS3}\\[8pt]
{\Large\color{steelblue} Metodología completa del último entrenamiento\\
split 361/63/84 · 100 épocas}\\[20pt]
\begin{tikzpicture}\draw[primaryblue, thick] (0,0) -- (11,0);\end{tikzpicture}

\vspace{14pt}
{\large \textit{Qué entra a cada red, qué le pasa a cada dato en cada capa,\\
de dónde salen los pesos, qué esperábamos y cómo se calcula cada número.}}

\vspace{22pt}
\includegraphics[width=0.95\textwidth]{ingesta_3d.png}

\vspace{14pt}
\begin{tcolorbox}[colback=panelbg,colframe=primaryblue,arc=4pt,width=0.92\textwidth,
  boxrule=1pt, left=14pt,right=14pt,top=8pt,bottom=8pt]
\begin{center}
\textbf{oasis3-ssl-workbench · \texttt{campaign\_ssl3d\_runpod/split361\_100}}\\[4pt]
\small 508 sujetos, split por sujeto 361/63/84 · VICReg, BYOL y DINOv2 con su receta campeona\\
100 épocas · 3 semillas · comparados con el supervisado completo, el encoder sin SSL\\
y el mismo encoder con pesos aleatorios
\end{center}
\end{tcolorbox}
\vfill
{\small 9 de octubre de 2026}
\end{titlepage}

\tableofcontents
\newpage

% ============================================================
\section{Cómo leer este documento}
% ============================================================

El documento sigue a un dato desde la resonancia hasta el número final. Primero se describe qué
volumen se recorta y cómo se convierte en un tensor (Secciones~\ref{sec:data}--\ref{sec:views}).
Después, de dónde salen los pesos de cada red y cómo se pasan de 2D a 3D
(Secciones~\ref{sec:weights}--\ref{sec:inflate}). Luego, qué hace cada capa con ese tensor
(Sección~\ref{sec:arch}), qué se compara (Sección~\ref{sec:conditions}) y cómo se calcula cada
métrica (Sección~\ref{sec:metrics}). Al final están los resultados esperados
(Sección~\ref{sec:expected}) y los obtenidos (Sección~\ref{sec:results}). No se explica la teoría
de cada método SSL; solo qué se calcula.

\begin{figure}[H]
\centering
\begin{tikzpicture}[x=1cm,y=1cm]
\node[stage] (mri) at (0,0) {\textbf{Resonancia T1}\\$97\times115\times97$ vóxeles\\de 2\,mm (MNI)};
\node[stage] (cube) at (3.1,0) {\textbf{Cubos fijos}\\$38^3$ guardados,\\$32^3$ evaluados};
\node[stage] (tensor) at (6.2,0) {\textbf{Tensor}\\$2\times32^3$ (CNN) o\\$6\times32^3$ (ViT)};
\node[stage] (views) at (9.3,0) {\textbf{Dos vistas}\\augmentation 3D\\de cada sujeto};
\node[stagew] (enc) at (9.3,-2.1) {\textbf{Encoder 3D}\\ConvNeXt, ResNet\\o ViT};
\node[stagew] (heads) at (6.2,-2.1) {\textbf{Cabezas}\\(solo para\\entrenar)};
\node[stageo] (loss) at (3.1,-2.1) {\textbf{Pérdida SSL}\\sin etiquetas};
\node[stage] (clean) at (0,-4.9) {\textbf{Cubo limpio}\\$32^3$ central, sin\\augmentation};
\node[stagew] (frozen) at (3.1,-4.9) {\textbf{Encoder}\\congelado, tras\\100 épocas};
\node[stage] (vec) at (6.2,-4.9) {\textbf{Un vector}\\por sujeto\\(768/512/384)};
\node[stageo] (probe) at (9.3,-4.9) {\textbf{Probe lineal}\\AD vs HC $\rightarrow$\\AUROC, F1, AUPR};
\draw[arr] (mri) -- (cube); \draw[arr] (cube) -- (tensor); \draw[arr] (tensor) -- (views);
\draw[arr] (views) -- (enc); \draw[arr] (enc) -- (heads); \draw[arr] (heads) -- (loss);
\draw[arr, dashed] (loss.south) -- ++(0,-0.32) -| node[pos=0.25, below, font=\scriptsize, text=primaryblue]
  {la pérdida actualiza encoder y cabezas} (enc.south);
\draw[arrg] (clean) -- (frozen); \draw[arrg] (frozen) -- (vec); \draw[arrg] (vec) -- (probe);
\node[font=\scriptsize\bfseries, text=primaryblue, anchor=west] at (-1.3,1.05) {Entrenamiento SSL (sin etiquetas)};
\node[font=\scriptsize\bfseries, text=softgray, anchor=west] at (-1.3,-3.9) {Evaluación (con etiquetas AD/HC)};
\end{tikzpicture}
\caption{El flujo completo. Arriba, el entrenamiento SSL: solo ve imágenes, nunca el
diagnóstico. Abajo, la evaluación: el encoder ya no cambia y una regresión logística (el
\emph{probe}) mide qué tan bien su vector separa Alzheimer (AD) de controles sanos (HC).}
\label{fig:flow}
\end{figure}

\begin{ideabox}[Vocabulario mínimo]
\begin{itemize}[leftmargin=12pt]
\item \textbf{Vóxel}: el píxel en 3D; aquí mide $2\times2\times2$\,mm.
\item \textbf{Cubo}: un trozo de la resonancia de $32\times32\times32$ vóxeles (64\,mm de lado).
\item \textbf{Canal}: una de las capas apiladas de la entrada, como R, G y B en una foto. Aquí
  cada canal es un cubo de una región distinta.
\item \textbf{Vista}: una copia del cubo transformada al azar (recortada, con ruido, deformada).
  Los métodos SSL aprenden comparando dos vistas del mismo sujeto.
\item \textbf{Encoder}: la red que convierte el cubo en un vector (la representación o
  \emph{embedding}). Es lo único que se conserva tras el SSL.
\item \textbf{Cabezas} (proyector, predictor, cabezas DINO): capas extra que solo existen para
  calcular la pérdida; se descartan al evaluar.
\item \textbf{Probe}: una regresión logística entrenada sobre los vectores del encoder congelado.
\item \textbf{Época}: una pasada por los 361 sujetos de entrenamiento. \textbf{Lote}: 16 sujetos.
  \textbf{Semilla}: el número que fija el azar; cada modelo se entrenó con 3 semillas.
\end{itemize}
\end{ideabox}

% ============================================================
\section{Qué se entrenó en la última campaña}
\label{sec:campaign}
% ============================================================

\begin{tldrbox}
Los tres modelos SSL campeones, 100 épocas y 3 semillas cada uno, más tres referencias con el mismo
encoder: entrenado de forma supervisada, sin SSL y con pesos aleatorios.
\end{tldrbox}

\begin{table}[H]
\centering
\small
\begin{tabularx}{\textwidth}{l>{\raggedright\arraybackslash}X>{\raggedright\arraybackslash}p{4.2cm}r}
\toprule
Corrida & Qué es & Nombre en disco & Semillas \\
\midrule
SSL \textcolor{vicregc}{\textbf{VICReg}} & ConvNeXt-T 3D, hipocampos $2\times32^3$, receta \texttt{geo}, validación con vistas aumentadas & \nolinkurl{vicreg_campaign_augval__geo__roi32} & 3 \\
SSL \textcolor{byolc}{\textbf{BYOL}} & ResNet-18 3D, hipocampos $2\times32^3$, receta \texttt{noise}, cabezas con GroupNorm & \nolinkurl{byol_gn_dp0.3_dq0_wd0.01__noise__roi32} & 3 \\
SSL \textcolor{dinoc}{\textbf{DINOv2}} & ViT-S 3D, 3 regiones $6\times32^3$, receta \texttt{elastic} + máscaras iBOT & \nolinkurl{dinov2_campaign__elastic__roimulti3} & 3 \\
\midrule
Supervisado & El mismo encoder + una capa lineal, entrenados con las etiquetas AD/HC & \nolinkurl{<modelo>_supervised__*} & 3 \\
Sin SSL & El encoder \emph{antes} del SSL (pesos ImageNet inflados), congelado & \nolinkurl{split361_30/<modelo>_init__*} & 1 \\
Aleatorio (ref.) & El mismo encoder con pesos aleatorios, congelado & \nolinkurl{split361_100/<modelo>_init__*} & 3 \\
\bottomrule
\end{tabularx}
\caption{Las corridas de \nolinkurl{campaign_ssl3d_runpod/split361_100} (más la corrida sin SSL, que
está en \texttt{split361\_30} con el mismo split y el mismo probe). RTX 5090 en RunPod, precisión
mixta bf16, 2.16 USD en total.}
\label{tab:runs}
\end{table}

% ============================================================
\section{Los datos de partida}
\label{sec:data}
% ============================================================

\begin{tldrbox}
Cada sesión es una resonancia T1 en espacio MNI de $97\times115\times97$ vóxeles de 2\,mm, con su
segmentación anatómica. Se normaliza a $[0,1]$ y se reparte por sujeto en 361/63/84.
\end{tldrbox}

\subsection{El volumen}

Los volúmenes vienen del LMDB de OASIS3 del proyecto (\texttt{data/dementia/mni-dataset.lmdb}). Cada
sesión trae:
\begin{itemize}
\item \textbf{\texttt{image}}: la resonancia T1 ya registrada al espacio MNI, en una grilla de
  $97\times115\times97$ vóxeles de 2\,mm, enteros de 16 bits. Los tres ejes del arreglo son
  $x$ (izquierda $\rightarrow$ derecha), $y$ (posterior $\rightarrow$ anterior) y $z$ (inferior
  $\rightarrow$ superior).
\item \textbf{\texttt{roi}}: la segmentación en la misma grilla, con las etiquetas estilo FreeSurfer
  (aseg + aparc): 17 y 53 son los hipocampos, 4 y 43 los ventrículos laterales, etc. Solo se usa
  para \emph{ubicar} los cubos y para medir volúmenes en el análisis del espacio latente; nunca
  entra a la red.
\end{itemize}

\textbf{Normalización de intensidad} (\texttt{normalize\_volume}, por volumen): se toman los
vóxeles distintos de cero, se calculan sus percentiles 1 y 99, y cada vóxel se transforma en
$(v - p_1)/(p_{99}-p_1)$, recortado a $[0,1]$. El fondo queda en 0. Así todas las resonancias
quedan en la misma escala, con el tejido entre 0 y 1.

\subsection{El split}

\begin{figure}[H]
\centering
\includegraphics[width=0.86\textwidth]{split_balance.png}
\caption{508 sujetos repartidos \emph{por sujeto}: ninguno aparece en dos conjuntos. Val y test solo
tienen AD y HC; ``otros'' (sin diagnóstico, Lewy, frontotemporal, vascular o etiqueta ambigua) solo
entran al SSL, que no usa etiquetas. Estratificado por diagnóstico y tercil de edad (semilla 42).}
\label{fig:split}
\end{figure}

\begin{table}[H]
\centering
\small
\begin{tabular}{lrrrrl}
\toprule
Split & Sujetos & HC & AD & Otros & Sesiones \\
\midrule
Train & 361 & 135 & 131 & 95 & 433 (varias por sujeto) \\
Val & 63 & 32 & 31 & 0 & 63 (una por sujeto) \\
Test & 84 & 43 & 41 & 0 & 84 (una por sujeto) \\
\bottomrule
\end{tabular}
\caption{Composición del split (\nolinkurl{campaign_ssl3d_runpod/split_361_63_84/SPLIT.md}).}
\end{table}

\begin{itemize}
\item \textbf{Entrenamiento SSL}: en cada época, cada uno de los 361 sujetos aporta \emph{una sesión
  elegida al azar} entre las suyas. Así una persona con 4 resonancias no pesa 4 veces más.
\item \textbf{Validación}: los 63 sujetos de val sirven para la pérdida de validación del SSL, para
  elegir la regularización del probe y para elegir la época del supervisado.
\item \textbf{Test}: los 84 sujetos de test solo se usan para el número final.
\end{itemize}

\begin{warnbox}[La edad viene pegada al diagnóstico]
En los tres conjuntos los AD tienen unos 7 años más que los HC (HC 67--69 años, AD 75--76). Parte de
lo que separa el probe puede ser envejecimiento y no Alzheimer (Sección~\ref{sec:latent}).
\end{warnbox}

% ============================================================
\section{Ingesta: de la resonancia al tensor}
\label{sec:ingest}
% ============================================================

\begin{tldrbox}
De cada resonancia se recortan cubos 3D en posiciones fijas; el del hemisferio derecho se espeja; los
cubos se apilan como canales. VICReg y BYOL reciben $2\times32\times32\times32$ (los dos
hipocampos) y DINOv2 $6\times32\times32\times32$ (hipocampo, ventrículo lateral y corteza
posteromedial, de cada lado). \textbf{Nada de esto es 2D}: la red recibe el volumen completo.
\end{tldrbox}

\begin{figure}[H]
\centering
\includegraphics[width=\textwidth]{ingesta_3d.png}
\caption{La ingesta en 3D. \textbf{A}: dónde están los seis cubos dentro del volumen MNI.
\textbf{B}: el cubo de $38^3$ que se guarda y el de $32^3$ que entra a la red. \textbf{C}: el
espejado del hemisferio derecho. \textbf{D}: los tensores de entrada; cada cubo está dibujado por sus
caras exteriores (lo que se ve desde fuera del volumen), no por un corte.}
\label{fig:ingesta}
\end{figure}

\subsection{Paso 1: dónde se corta cada cubo}

La posición de cada cubo es \textbf{la misma para todos los sujetos}: el centroide medio de la
región, calculado una sola vez sobre la segmentación de las sesiones de entrenamiento (429 sesiones
con segmentación; solo anatomía, nunca el diagnóstico). Con el train del split 361 los centroides
redondeados dan exactamente las mismas posiciones, por eso la caché se reutilizó. Como las imágenes
están registradas a MNI, una posición fija cae sobre la misma estructura en todos los cerebros, con
la variación anatómica natural de cada persona.

Alrededor de cada centroide se recorta un cubo de \textbf{38 vóxeles} de lado
($38 = \mathrm{round}(1.2\times32)$), que empieza en $\mathrm{round}(\text{centro}) - 19$ en cada
eje. Ese es el cubo que se guarda en disco (\texttt{data/roi3d/hippo\_s32.npy} y
\texttt{multi3\_s32.npy}, en float16, 594 sesiones). La red solo ve los \textbf{32 vóxeles
centrales} (3 de margen por lado): el margen sirve para que los recortes y deformaciones de las
vistas no salgan del cubo.

\begin{table}[H]
\centering
\small
\begin{tabular}{llrrrr}
\toprule
Región & Lado & Centroide $(x,y,z)$ & \multicolumn{3}{c}{Cubo evaluado de $32^3$ (índices de vóxel)} \\
\cmidrule(lr){4-6}
 & & & $x$ & $y$ & $z$ \\
\midrule
\multirow{2}{*}{Hipocampo} & izq. & (35.4, 54.9, 32.2) & 19--50 & 39--70 & 16--47 \\
 & der. & (61.2, 55.4, 32.3) & 45--76 & 39--70 & 16--47 \\
\multirow{2}{*}{Ventrículo lateral} & izq. & (41.3, 58.8, 46.1) & 25--56 & 43--74 & 30--61 \\
 & der. & (55.2, 59.9, 46.1) & 39--70 & 44--75 & 30--61 \\
\multirow{2}{*}{Corteza posteromedial} & izq. & (44.5, 42.0, 56.9) & 29--60 & 26--57 & 41--72 \\
 & der. & (52.4, 42.4, 57.1) & 36--67 & 26--57 & 41--72 \\
\bottomrule
\end{tabular}
\caption{Posiciones exactas (\texttt{data/roi3d/multi3\_s32.json}). La corteza posteromedial es el
cíngulo posterior e ístmico más el precúneo (etiquetas 1023, 1010, 1025 y sus pares 2xxx). Ojo:
esa región y el ventrículo están junto a la línea media, así que sus cubos izquierdo y derecho se
solapan mucho (25 y 18 de 32 vóxeles en $x$); los del hipocampo apenas (6).}
\label{tab:positions}
\end{table}

\subsection{Paso 2: el hemisferio derecho se espeja}

El cubo derecho se voltea en el eje $x$ (\texttt{np.flip(cube, axis=0)}) antes de guardarlo
(Figura~\ref{fig:ingesta}\,C). Así el hipocampo derecho queda orientado igual que el izquierdo, y un
mismo filtro sirve para ambos. Consecuencia útil: voltear la cabeza de izquierda a derecha equivale
exactamente a \emph{intercambiar los dos canales de cada par}, que es como se implementa esa
augmentation.

\subsection{Paso 3: los cubos se apilan como canales}

Los cubos \textbf{no} se pegan uno junto a otro ni se mezclan: se apilan en una dimensión nueva, la
de canales, igual que R, G y B en una foto.

\begin{table}[H]
\centering
\small
\begin{tabularx}{\textwidth}{l>{\raggedright\arraybackslash}Xll}
\toprule
Modelo & Canales, en orden & Tensor por sujeto & Lote de 16 \\
\midrule
VICReg y BYOL & 0: hipocampo izq.; 1: hipocampo der. espejado & $2\times32\times32\times32$ & $16\times2\times32\times32\times32$ \\
DINOv2 & 0: hipocampo izq.; 1: hipocampo der. esp.; 2: ventrículo izq.; 3: ventrículo der. esp.;
  4: posteromedial izq.; 5: posteromedial der. esp. & $6\times32\times32\times32$ & $16\times6\times32\times32\times32$ \\
\bottomrule
\end{tabularx}
\caption{Los tensores de entrada. Cada número es la intensidad normalizada ($[0,1]$) de un vóxel.}
\label{tab:tensors}
\end{table}

\textbf{Cómo DINOv2 llega a 6 canales.} La caché \texttt{multi3\_s32} guarda, por sesión, un
arreglo de $6\times38\times38\times38$: para cada una de las 3 regiones, el cubo izquierdo y el
derecho espejado (\texttt{build\_region\_cache}). La primera capa del ViT simplemente acepta 6
canales de entrada en vez de 3 (Sección~\ref{sec:inflate}). Un detalle importante: como los 6
cubos se apilan, la posición $(i,j,k)$ de cada canal corresponde a un lugar anatómico
\emph{distinto} en cada región. Un parche del ViT mezcla, en un solo token, los 6 cubos en la misma
posición relativa (Figura~\ref{fig:volumes}). La red tiene que aprender qué canal es qué región.

\begin{warnbox}[2D o 3D: la confusión de la presentación]
Todo el pipeline es 3D. Las figuras con cortes (axial, coronal, sagital) son solo una forma de
\emph{mostrar} un volumen en papel; la red nunca ve cortes sueltos. Lo único 2D en toda la historia
son los \textbf{pesos de partida}: vienen de redes entrenadas con fotos 2D de ImageNet y se
convierten a 3D antes de empezar (Sección~\ref{sec:inflate}). Cuando una figura dice
``parche $16^2\rightarrow8^3$'' significa eso: el ViT original usaba parches de $16\times16$
píxeles y el nuestro usa parches cúbicos de $8\times8\times8$ vóxeles.
\end{warnbox}

% ============================================================
\section{Las vistas: qué augmentation recibe cada modelo}
\label{sec:views}
% ============================================================

\begin{tldrbox}
Cada sujeto del lote produce dos vistas al azar con la receta de su modelo. Todas las operaciones
son 3D y actúan sobre el cubo de $38^3$ antes de quedarse con $32^3$. El probe, en cambio, siempre
recibe el cubo central limpio.
\end{tldrbox}

\subsection{El orden de las operaciones}

\texttt{ROIView3DTransform} genera cada vista en este orden (todas las operaciones existen; cada
receta activa solo algunas):
\begin{enumerate}
\item \textbf{Artefactos de TorchIO} sobre el cubo completo de $38^3$ y sus $C$ canales a la vez
  (ruido, deformación elástica\ldots).
\item \textbf{Recorte}: un sub-cubo de lado $s$ vóxeles, en el centro o en una posición al azar.
\item \textbf{Reescalado} trilineal del sub-cubo a $32^3$ (si $s=32$, no cambia nada).
\item \textbf{Volteo izquierda--derecha} con probabilidad 0.5: intercambio de los canales de cada par.
\item \textbf{Rotación} 3D: un ángulo al azar en cada eje, aplicada con \texttt{grid\_sample}
  (interpolación trilineal, relleno con el borde).
\item Cambios fotométricos (desenfoque, escala de intensidad): ninguna receta campeona los usa.
\item \textbf{Recorte a $[0,1]$} de todos los vóxeles.
\end{enumerate}

\begin{table}[H]
\centering
\small
\begin{tabularx}{\textwidth}{l>{\raggedright\arraybackslash}X>{\raggedright\arraybackslash}X>{\raggedright\arraybackslash}X}
\toprule
 & \textcolor{vicregc}{\textbf{VICReg · \texttt{geo}}} & \textcolor{byolc}{\textbf{BYOL · \texttt{noise}}} & \textcolor{dinoc}{\textbf{DINOv2 · \texttt{elastic}}} \\
\midrule
1. TorchIO & no & \texttt{RandomNoise}: $\sigma\sim U(0,\,0.215)$ por vista, ruido gaussiano independiente en cada vóxel y canal & \nolinkurl{RandomElasticDeformation}: rejilla de $7\times7\times7$ puntos de control, desplazamiento $\leq1.95$ vóxeles (3.9\,mm), bordes fijos; el mismo campo en los 6 canales \\
2--3. Recorte & sub-cubo de 32 a 35 vóxeles, posición al azar, reescalado a 32 & cubo central de 32 (sin recorte) & cubo central de 32 (sin recorte) \\
4. Volteo izq.--der. & $p=0.5$ & $p=0.5$ & $p=0.5$ (los 3 pares) \\
5. Rotación & $\pm5^\circ$ en cada eje & no & no \\
Además & -- & -- & máscara iBOT sobre los parches, dentro de la red (Sección~\ref{sec:vit}) \\
\bottomrule
\end{tabularx}
\caption{Las recetas campeonas (\texttt{t5\_ssl.recipes.AUG\_RECIPES}). La intensidad del ruido y de
la deformación está calibrada para alterar el cubo lo mismo (RMSE 0.08). Como el ruido se suma antes
de recortar a $[0,1]$, en el fondo (que vale 0) solo sobrevive el ruido positivo.}
\label{tab:aug}
\end{table}

\begin{figure}[H]
\centering
\includegraphics[width=0.78\textwidth]{views.png}
\caption{Cubo limpio (lo que ve el probe) y dos vistas de cada receta, para el canal del hipocampo
izquierdo de un sujeto de validación. Cada imagen apila los cortes centrales de los tres planos
\emph{del mismo volumen}: solo es una forma de dibujarlo.}
\label{fig:views}
\end{figure}

\subsection{Qué vistas se usan en cada momento}

\begin{table}[H]
\centering
\small
\begin{tabularx}{\textwidth}{l>{\raggedright\arraybackslash}X}
\toprule
Momento & Entrada \\
\midrule
Entrenamiento SSL & Dos vistas al azar de la receta, nuevas en cada época. \\
Pérdida de validación SSL & VICReg: dos vistas de la receta, fijas por sesión (misma semilla en cada
  época, \texttt{val\_views=train}). BYOL y DINOv2: el cubo limpio y su espejo izquierda--derecha. \\
Probe, monitor de colapso, espacio latente & El cubo central de $32^3$, sin ninguna augmentation ni
  volteo (\texttt{ROIEval3D}). \\
Supervisado (entrenamiento) & Una sola vista de la misma receta del modelo. \\
\bottomrule
\end{tabularx}
\end{table}

\subsection{Lotes y pasos}

Un lote son 16 sujetos. Con 361 sujetos hay 23 lotes por época (22 de 16 y uno de 9): 2300 pasos de
optimización en 100 épocas. En cada paso las redes procesan:
\begin{itemize}
\item \textbf{VICReg}: las dos vistas por el mismo encoder, $2\times16 = 32$ volúmenes.
\item \textbf{BYOL}: las dos vistas por la red \emph{online} y las dos por la copia \emph{target}:
  64 pasadas de volumen, de las cuales solo las 32 de la online reciben gradiente.
\item \textbf{DINOv2}: las dos vistas enmascaradas por el \emph{alumno} y las dos limpias por el
  \emph{profesor}: también 64, con gradiente solo en el alumno.
\end{itemize}

% ============================================================
\section{Los pesos de partida: de dónde salen}
\label{sec:weights}
% ============================================================

\begin{tldrbox}
Los tres encoders arrancan de pesos públicos entrenados con \textbf{fotos naturales de ImageNet-1k
(clasificación supervisada en 2D)}. Ninguno se entrenó con resonancias ni con imágenes médicas. Se
convierten a 3D (Sección~\ref{sec:inflate}) y ese es el punto de partida del SSL, del supervisado y
de la barra ``sin SSL''.
\end{tldrbox}

\begin{table}[H]
\centering
\small
\begin{tabularx}{\textwidth}{l>{\raggedright\arraybackslash}p{3.2cm}>{\raggedright\arraybackslash}X>{\raggedright\arraybackslash}p{3.6cm}}
\toprule
Encoder & Pesos & Origen exacto & Con qué se entrenaron \\
\midrule
\textcolor{vicregc}{\textbf{ConvNeXt-T}} (VICReg) & torchvision \nolinkurl{ConvNeXt_Tiny_Weights.IMAGENET1K_V1} &
  \nolinkurl{https://download.pytorch.org/models/convnext_tiny-983f1562.pth} (Liu et al., 2022, \emph{A ConvNet for the 2020s}) &
  ImageNet-1k, 1.28\,M fotos RGB, 1000 clases; 82.5\,\% top-1 \\
\textcolor{byolc}{\textbf{ResNet-18}} (BYOL) & torchvision \nolinkurl{ResNet18_Weights.IMAGENET1K_V1} &
  \nolinkurl{https://download.pytorch.org/models/resnet18-f37072fd.pth} (He et al., 2016) &
  ImageNet-1k; 69.8\,\% top-1 \\
\textcolor{dinoc}{\textbf{ViT-S/16}} (DINOv2) & DeiT-S (\nolinkurl{deit_small_patch16_224}) &
  \nolinkurl{https://dl.fbaipublicfiles.com/deit/deit_small_patch16_224-cd65a155.pth} (Touvron et al., 2021, \emph{Training data-efficient image transformers}) &
  ImageNet-1k, supervisado, sin destilación; $\approx$80\,\% top-1 \\
\bottomrule
\end{tabularx}
\caption{Los pesos usados. Se descargan una vez a \texttt{\textasciitilde/.cache/torch/hub/checkpoints/}.
En un host de RunPod que no llegaba a \texttt{download.pytorch.org}, la ResNet-18 se bajó de Hugging
Face (\texttt{timm/resnet18.tv\_in1k}), que contiene los mismos tensores.}
\label{tab:weights}
\end{table}

\begin{itemize}
\item \textbf{Por qué DeiT-S y no los pesos oficiales de DINOv2.} El código admite también
  \texttt{dinov2\_vits14} (DINOv2 ViT-S/14, autosupervisado en 142 millones de imágenes
  naturales), pero la campaña usa DeiT-S para que los tres modelos partan de lo mismo: una red
  \emph{supervisada en ImageNet-1k}. Así la comparación entre métodos no mezcla puntos de partida.
\item \textbf{Las cabezas} (proyector, predictor, cabezas DINO) no tienen pesos preentrenados: se
  inicializan al azar.
\item \textbf{La referencia ``aleatoria''} usa la arquitectura idéntica con la inicialización por
  defecto de PyTorch (CNN) o normal truncada con $\sigma=0.02$ (ViT), sin ningún peso de ImageNet.
\end{itemize}

% ============================================================
\section{De 2D a 3D: qué se hace exactamente con los pesos}
\label{sec:inflate}
% ============================================================

\begin{tldrbox}
Cada filtro 2D de $k\times k$ se repite $k$ veces en profundidad y se divide por $k$ (``inflado''
de I3D). Los canales RGB se colapsan en uno gris y se reparten entre los 2 o 6 canales de entrada.
En el ViT, además, el parche de $16\times16$ se reduce a $8\times8$ (redimensionado de FlexiViT) y
las posiciones aprendidas se extienden a una rejilla de $4\times4\times4$.
\end{tldrbox}

\subsection{Las CNN: inflado I3D}

\textbf{Base}: Carreira \& Zisserman (2017), \emph{Quo Vadis, Action Recognition? A New Model and the
Kinetics Dataset} (CVPR), donde se ``inflan'' redes de ImageNet para video. Código:
\texttt{src/t5\_ssl/models/backbones/inflate3d.py}.

\begin{enumerate}
\item \textbf{Cada convolución} \texttt{Conv2d} de $k\times k$ pasa a \texttt{Conv3d} de
  $k\times k\times k$ con el mismo paso y relleno:
  \[ W_{3D}[o,i,d,h,w] = \frac{W_{2D}[o,i,h,w]}{k}, \qquad d = 1,\ldots,k. \]
  Propiedad: un volumen que es constante a lo largo de la profundidad produce, en cada corte, las
  mismas activaciones que la red 2D sobre ese corte (salvo en los bordes, por el relleno). La red
  empieza ``sabiendo'' lo mismo que en 2D y el entrenamiento la vuelve verdaderamente 3D.
\item \textbf{La primera capa} pasa de 3 canales (RGB) a $C$ canales: se suman los tres filtros
  RGB (la respuesta a una imagen gris) y el resultado se copia en cada uno de los $C$ canales,
  dividido por $C$:
  \[ W_{3D}^{\,\text{tallo}}[o,c,d,h,w] = \frac{1}{C\,k}\sum_{\text{RGB}} W_{2D}[o,\text{rgb},h,w]. \]
\item \textbf{Normalizaciones}: \texttt{BatchNorm2d} $\rightarrow$ \texttt{BatchNorm3d} y
  \texttt{LayerNorm2d} $\rightarrow$ \texttt{LayerNorm3d}, copiando sus parámetros y estadísticas.
  El escalado por canal de los bloques ConvNeXt (\emph{layer scale}) se copia con una dimensión
  más; el promedio global pasa a 3D; se descarta el clasificador de 1000 clases.
\item \textbf{ResNet-18}: se elimina el \emph{max-pool} del tallo. Con cubos de 32 vóxeles, tallo de
  paso 2 más max-pool dejarían la última capa en $1^3$.
\end{enumerate}

\begin{ideabox}[Tres consecuencias que importan]
\begin{itemize}[leftmargin=12pt]
\item \textbf{Qué plano ven los filtros 2D.} La profundidad del filtro es el primer eje del cubo, $x$
  (izquierda--derecha). Los filtros de ImageNet quedan aplicados sobre planos \emph{sagitales}
  ($y$--$z$) y promediados a lo largo de $x$.
\item \textbf{Al inicio, la red ve el promedio de los canales.} Todos los canales comparten el
  mismo filtro, así que la primera capa responde a la suma de los cubos. Intercambiar canales (el
  volteo izquierda--derecha) no cambia nada al inicio: el SSL tiene que romper esa simetría.
\item \textbf{Sin normalización de ImageNet en las CNN.} La entrada es $[0,1]$ sin restar la media de
  ImageNet; la BatchNorm (ResNet) o la LayerNorm (ConvNeXt) tras el tallo reescala las activaciones.
\end{itemize}
\end{ideabox}

\subsection{El ViT: inflado del parche y de las posiciones}
\label{sec:vitinflate}

Código: \texttt{inflate\_patch\_embed} e \texttt{inflate\_pos\_embed} en
\nolinkurl{src/t5_ssl/models/backbones/vit3d.py}. Son cinco pasos sobre los pesos de DeiT-S:

\begin{enumerate}
\item \textbf{Normalización de ImageNet absorbida.} DeiT espera cada canal RGB normalizado como
  $(x-\mu_c)/\sigma_c$. Como nuestra entrada es gris en $[0,1]$, esa resta y división se pasan al
  filtro: $W_\text{gris} = \sum_c W_c/\sigma_c$ y $b' = b - \sum_c \sum_{h,w} W_c\,\mu_c/\sigma_c$.
  Es álgebra exacta: el token de un parche gris es el mismo que DeiT daría a ese parche en RGB.
\item \textbf{De parche $16\times16$ a $8\times8$.} Se resuelve qué filtro de $8\times8$ da, para
  cada parche, el mismo token que el filtro original da a ese parche \emph{ampliado} a $16\times16$
  por interpolación bilineal: $W_8 = U^{\top} W_{16}$, donde $U$ es la matriz de esa ampliación
  ($256\times64$). Equivale a mostrarle a DeiT la imagen con zoom $\times2$. La idea de
  redimensionar el filtro con el operador de interpolación viene de FlexiViT (Beyer et al., 2023,
  \emph{FlexiViT: One Model for All Patch Sizes}, CVPR).
\item \textbf{De $8\times8$ a $8\times8\times8$ y de 1 a 6 canales.} Como en I3D: el filtro se repite
  8 veces en profundidad y en los 6 canales, y se divide por $8\times6=48$. Es la inicialización por
  ``inflado de filtros'' del \emph{tubelet embedding} de ViViT (Arnab et al., 2021, \emph{ViViT: A
  Video Vision Transformer}, ICCV).
\item \textbf{Posiciones.} DeiT tiene 196 posiciones aprendidas (una rejilla de $14\times14$) más la
  del CLS. La rejilla se reduce a $4\times4$ por interpolación bicúbica (lo habitual al cambiar la
  resolución de un ViT, Dosovitskiy et al., 2021), se repite en las 4 capas de profundidad (como en
  ViViT) y a cada capa $d$ se le suma el ``código'' centrado de la fila $d$ (el promedio de esa fila
  de la rejilla). Este último ajuste es propio del repositorio: sirve para que las 64 posiciones 3D
  empiecen distintas entre sí. La posición del CLS se copia.
\item \textbf{Todo lo demás se copia sin cambios}: los 12 bloques Transformer, la LayerNorm final y
  el token CLS. Esos bloques operan sobre una secuencia de tokens y no saben si vienen de 2D o 3D. El
  \emph{mask token} no existe en DeiT: se inicializa al azar (depende de la semilla) y solo se usa
  para tapar parches durante el entrenamiento.
\end{enumerate}

\begin{warnbox}[Por qué el ViT sin SSL parte tan bajo]
Con el filtro compartido, al inicio el token de un parche es el token DeiT del \emph{promedio} de sus
6 canales y de sus 8 cortes sagitales. Es decir: el ViT de partida mezcla hipocampo, ventrículo y
corteza posteromedial en una sola imagen. Es una explicación probable (no demostrada) de que ese
encoder sea el peor punto de partida (AUROC 0.778 en CV y 0.727 en test, por debajo incluso del ViT
con pesos aleatorios) y de que DINOv2 sea el que más gana con el SSL.
\end{warnbox}

% ============================================================
\section{Las arquitecturas, capa por capa}
\label{sec:arch}
% ============================================================

\begin{tldrbox}
Las tres redes mantienen una rejilla espacial 3D hasta el final y la colapsan en un vector con un
promedio (CNN) o con el token CLS (ViT). Ese vector (768, 512 o 384 números) es lo que mide el
probe. Las cabezas y la copia lenta (EMA) solo existen para entrenar.
\end{tldrbox}

\begin{figure}[H]
\centering
\includegraphics[width=\textwidth]{volumenes_3d.png}
\caption{Qué le pasa al volumen dentro de cada encoder. En las CNN, cada caja es un mapa de
activaciones 3D: el lado indica el tamaño de la rejilla ($8^3$, $4^3$\ldots) y la pila, los canales.
La caja de color sobre la entrada es la ventana que ve una posición del primer mapa. En el ViT, el
cubo de $6\times32^3$ se corta en $4\times4\times4=64$ parches cúbicos de $8^3$; cada parche (con sus
6 canales) se convierte en un token de 384 números. El bloque violeta es un ejemplo de máscara iBOT.}
\label{fig:volumes}
\end{figure}

Las tablas siguientes dan la forma del tensor después de cada capa \textbf{para un sujeto}; con un
lote se antepone la dimensión 16. ``Parám.'' es el número de pesos de esa parte.

\subsection{VICReg: ConvNeXt-T 3D}
\label{sec:convnext}

\begin{table}[H]
\centering
\small
\begin{tabularx}{\textwidth}{l>{\raggedright\arraybackslash}Xlr}
\toprule
Parte & Operación & Salida & Parám. \\
\midrule
Entrada & cubo de 2 canales & $2\times32\times32\times32$ & -- \\
Tallo & conv.\ $4\times4\times4$ con paso 4 (parches 3D no solapados) + LayerNorm & $96\times8\times8\times8$ & 12.6\,k \\
Etapa 1 & 3 bloques ConvNeXt de 96 canales & $96\times8^3$ & 0.32\,M \\
Reducción & LayerNorm + conv.\ $2\times2\times2$ con paso 2 & $192\times4^3$ & 0.15\,M \\
Etapa 2 & 3 bloques de 192 & $192\times4^3$ & 1.09\,M \\
Reducción & LayerNorm + conv.\ $2^3$, paso 2 & $384\times2^3$ & 0.59\,M \\
Etapa 3 & 9 bloques de 384 & $384\times2^3$ & 11.83\,M \\
Reducción & LayerNorm + conv.\ $2^3$, paso 2 & $768\times1^3$ & 2.36\,M \\
Etapa 4 & 3 bloques de 768 & $768\times1^3$ & 14.97\,M \\
Promedio global & media sobre la rejilla & \textbf{768} & -- \\
\midrule
\multicolumn{3}{l}{\textbf{Total del encoder}} & \textbf{31.3\,M} \\
\bottomrule
\end{tabularx}
\caption{ConvNeXt-T 3D (\texttt{ConvNeXt3DBackbone}).}
\end{table}

\textbf{Un bloque ConvNeXt 3D} con $C$ canales hace, en orden:
\begin{enumerate}
\item convolución \emph{por canal} (depthwise) de $7\times7\times7$ con relleno 3: cada canal se
  filtra por separado, mirando un vecindario de 7 vóxeles en cada eje;
\item LayerNorm sobre los $C$ canales de cada posición;
\item lineal $C\rightarrow4C$, GELU, lineal $4C\rightarrow C$ (aplicadas posición por posición: mezclan
  canales, no posiciones);
\item multiplicación por un vector aprendido $\gamma$ (\emph{layer scale});
\item \emph{drop path}: en entrenamiento, con probabilidad que crece de 0 a 0.1 a lo largo de los 18
  bloques, se apaga el bloque entero para ese sujeto;
\item suma con la entrada del bloque (conexión residual).
\end{enumerate}
En la etapa 4 la rejilla es $1^3$: el filtro $7^3$ solo usa su peso central.

\textbf{Proyector} (solo entrenamiento, 10.0\,M): lineal $768\rightarrow2048$, BatchNorm, ReLU,
dropout 0.1, lineal $2048\rightarrow2048$, BatchNorm, ReLU, dropout 0.1, lineal
$2048\rightarrow2048$.

\textbf{Un paso de entrenamiento.} Las vistas $v_1$ y $v_2$ ($16\times2\times32^3$ cada una) pasan
por \emph{el mismo} encoder (mismos pesos) y dan $h_1,h_2$ ($16\times768$); el proyector da
$z_1,z_2$ ($16\times2048$). La pérdida, en precisión simple, suma
$25\cdot\mathrm{MSE}(z_1,z_2)$ (las dos vistas deben dar lo mismo), $25\cdot$ la penalización de
cada una de las 2048 dimensiones cuya desviación estándar en el lote de 16 baje de 1, y $1\cdot$ la
suma de covarianzas al cuadrado fuera de la diagonal (matriz de $2048\times2048$), dividida por 2048.
El probe usa $h$ (768).

\subsection{BYOL: ResNet-18 3D}
\label{sec:resnet}

\begin{table}[H]
\centering
\small
\begin{tabularx}{\textwidth}{l>{\raggedright\arraybackslash}Xlr}
\toprule
Parte & Operación & Salida & Parám. \\
\midrule
Entrada & cubo de 2 canales & $2\times32\times32\times32$ & -- \\
Tallo & conv.\ $7\times7\times7$, paso 2, relleno 3 + BatchNorm + ReLU (sin max-pool) & $64\times16\times16\times16$ & 44\,k \\
Capa 1 & 2 bloques básicos de 64 & $64\times16^3$ & 0.44\,M \\
Capa 2 & 2 bloques de 128; el primero con paso 2 & $128\times8^3$ & 1.56\,M \\
Capa 3 & 2 bloques de 256; el primero con paso 2 & $256\times4^3$ & 6.23\,M \\
Capa 4 & 2 bloques de 512; el primero con paso 2 & $512\times2^3$ & 24.91\,M \\
Promedio global & media sobre la rejilla & \textbf{512} & -- \\
\midrule
\multicolumn{3}{l}{\textbf{Total del encoder}} & \textbf{33.2\,M} \\
\bottomrule
\end{tabularx}
\caption{ResNet-18 3D (\texttt{ResNet3DBackbone}).}
\end{table}

\textbf{Un bloque básico 3D}: conv.\ $3\times3\times3$ $\rightarrow$ BatchNorm $\rightarrow$ ReLU
$\rightarrow$ conv.\ $3\times3\times3$ $\rightarrow$ BatchNorm $\rightarrow$ suma con el atajo
$\rightarrow$ ReLU. Cuando el bloque reduce la rejilla (paso 2) y cambia el número de canales, el
atajo es una conv.\ $1\times1\times1$ con paso 2 + BatchNorm.

\textbf{Cabezas} (solo entrenamiento), con GroupNorm (32 grupos de 128 neuronas) en vez de BatchNorm:
\begin{itemize}
\item \textbf{Proyector} (3.2\,M): lineal $512\rightarrow4096$, GroupNorm, ReLU, dropout 0.3, lineal
  $4096\rightarrow256$.
\item \textbf{Predictor} (2.1\,M, solo en la red online): lineal $256\rightarrow4096$, GroupNorm,
  ReLU, lineal $4096\rightarrow256$.
\end{itemize}

\textbf{La copia lenta (target).} Es una copia del encoder y del proyector (sin predictor, sin
dropout) que no recibe gradiente. Tras cada paso se actualiza como
$\theta_\text{target}\leftarrow\tau\,\theta_\text{target}+(1-\tau)\,\theta_\text{online}$, con
$\tau_t = 1-(1-0.996)\,\frac{\cos(\pi t/T)+1}{2}$: empieza en 0.996 y sube a 1 en los $T=2300$ pasos.
Las estadísticas de BatchNorm se copian directamente de la online.

\textbf{Un paso de entrenamiento.} La online procesa $v_1$ y $v_2$ hasta el predictor ($p_1,p_2$, de
$16\times256$); la target procesa las mismas vistas hasta su proyector ($\bar z_1,\bar z_2$). La
pérdida es $\tfrac12\big[(2-2\cos(p_1,\bar z_2))+(2-2\cos(p_2,\bar z_1))\big]$, promediada en el lote:
cada vista predice lo que la copia lenta dice de la \emph{otra}. El probe usa el encoder online (512).

\subsection{DINOv2: ViT-S 3D}
\label{sec:vit}

\begin{table}[H]
\centering
\small
\begin{tabularx}{\textwidth}{l>{\raggedright\arraybackslash}Xlr}
\toprule
Parte & Operación & Salida & Parám. \\
\midrule
Entrada & cubo de 6 canales & $6\times32\times32\times32$ & -- \\
Parcheado & conv.\ $8\times8\times8$ con paso 8: cada parche de $6\times8^3 = 3072$ valores $\rightarrow$ 384 & $384\times4\times4\times4$ & 1.18\,M \\
Aplanado & la rejilla $4\times4\times4$ se lee en orden $(i,j,k)$ & $64\times384$ & -- \\
Máscara (solo alumno) & los tokens de los parches enmascarados se reemplazan por el \emph{mask token} & $64\times384$ & 384 \\
CLS & se antepone un token aprendido & $65\times384$ & 384 \\
Posición & se suma una posición aprendida por token & $65\times384$ & 25\,k \\
12 bloques & Transformer (ver abajo) & $65\times384$ & 21.3\,M \\
LayerNorm final & & $65\times384$ & 768 \\
Salidas & token 0 = \textbf{CLS} y 64 tokens de parche & \textbf{384} y $64\times384$ & -- \\
\midrule
\multicolumn{3}{l}{\textbf{Total del encoder}} & \textbf{22.5\,M} \\
\bottomrule
\end{tabularx}
\caption{ViT-S 3D (\texttt{ViT3DBackbone}), parches de $8^3$.}
\end{table}

\textbf{Un bloque Transformer} (1.77\,M): LayerNorm $\rightarrow$ atención multi-cabeza (6 cabezas de
64 dimensiones; una lineal $384\rightarrow1152$ da consultas, claves y valores; cada token pondera a
los 65 tokens con $\mathrm{softmax}(QK^\top/\sqrt{64})$; lineal de salida $384\rightarrow384$)
$\rightarrow$ drop path $\rightarrow$ suma residual; luego LayerNorm $\rightarrow$ MLP
$384\rightarrow1536\rightarrow384$ con GELU $\rightarrow$ drop path $\rightarrow$ suma residual. El
drop path crece de 0 a 0.1 a lo largo de los 12 bloques. Gracias a la atención, cada parche puede
combinar información de \emph{cualquier} otro parche del cubo desde el primer bloque.

\textbf{La máscara iBOT} (\texttt{sample\_masks}): en cada lote y cada vista, 8 de las 16 imágenes
llevan máscara. A la imagen $n$-ésima de esas 8 le toca una fracción al azar dentro del tramo
$[0.1+0.4\,n/8,\ 0.1+0.4\,(n+1)/8]$, es decir, entre 6 y 31 de los 64 parches. Los parches se tapan
en bloques 3D contiguos de al menos 4 parches, con proporciones entre 0.3 y 3.3. El token tapado se
reemplaza \emph{antes} de sumar la posición: la red sabe dónde está el hueco, no qué hay en él.

\textbf{Cabezas} (solo entrenamiento, 6.6\,M cada una, con pesos separados):
\begin{itemize}
\item \textbf{Cabeza DINO} (sobre el CLS): lineal $384\rightarrow2048$, BatchNorm, GELU, lineal
  $2048\rightarrow2048$, BatchNorm, GELU, lineal $2048\rightarrow256$; el vector de 256 se normaliza
  (largo 1) y se compara por coseno con \textbf{4096 prototipos} (vectores aprendidos de largo 1):
  salen 4096 puntajes entre $-1$ y 1.
\item \textbf{Cabeza iBOT} (sobre cada token de parche enmascarado): la misma forma.
\end{itemize}

\textbf{El profesor.} Copia del encoder y de las dos cabezas, sin drop path ni gradiente, que se
actualiza como en BYOL con $\tau$ de 0.99 a 1.

\textbf{Un paso de entrenamiento.}
\begin{enumerate}
\item Se sortean las máscaras de las dos vistas.
\item \textbf{Alumno}: cada vista enmascarada $\rightarrow$ encoder $\rightarrow$ CLS $\rightarrow$
  cabeza DINO ($16\times4096$); los tokens de los parches tapados $\rightarrow$ cabeza iBOT.
\item \textbf{Profesor}: cada vista \emph{sin} máscara $\rightarrow$ encoder $\rightarrow$ CLS
  $\rightarrow$ su cabeza DINO; los tokens en las posiciones tapadas $\rightarrow$ su cabeza iBOT.
\item Los puntajes del profesor se convierten en probabilidades con Sinkhorn--Knopp (3 iteraciones,
  temperatura 0.01), que reparte el lote de forma pareja entre los 4096 prototipos; los del alumno,
  con un softmax a temperatura 0.1.
\item \textbf{DINO}: entropía cruzada entre el CLS del alumno en una vista y el del profesor en la
  \emph{otra} (los dos cruces, promediados). \textbf{iBOT}: entropía cruzada entre cada parche tapado
  del alumno y el mismo parche del profesor, con el mismo peso por imagen. \textbf{KoLeo}: castiga
  que dos CLS del alumno (normalizados) del mismo lote queden demasiado cerca.
\item Pérdida $=$ DINO $+\,1\cdot$iBOT $+\,0.1\cdot$KoLeo; gradiente recortado a norma 3; paso de
  AdamW; actualización EMA del profesor. Durante la primera época los prototipos no se actualizan.
\end{enumerate}
El probe usa el \textbf{CLS del profesor} (384), como en DINO y DINOv2.

\subsection{Las figuras de arquitectura}

\begin{figure}[H]
\centering
\includegraphics[width=\textwidth]{architectures_100.png}
\caption{Los tres campeones de la entrada a la pérdida, con la configuración de este entrenamiento
(100 épocas, BYOL con cabezas GroupNorm). Es la figura \texttt{architectures.png} del informe
anterior, que mostraba la de 30 épocas. El vector marcado \emph{probe} es lo que se evalúa; la caja
punteada es la copia lenta (EMA).}
\label{fig:architectures}
\end{figure}

\begin{figure}[H]
\centering
\includegraphics[width=\textwidth]{arch_vicreg.png}
\caption{VICReg rama por rama: dos vistas \texttt{geo} de un sujeto de validación por la misma
ConvNeXt-T 3D (pesos compartidos).}
\label{fig:archvicreg}
\end{figure}

\begin{figure}[H]
\centering
\includegraphics[width=\textwidth]{arch_byol.png}
\caption{BYOL: solo la rama online tiene predictor y recibe gradiente; la target (punteada) es la
copia EMA (\emph{sg} = sin gradiente).}
\label{fig:archbyol}
\end{figure}

\begin{figure}[H]
\centering
\includegraphics[width=\textwidth]{arch_dinov2.png}
\caption{DINOv2: el alumno ve la vista con un bloque de parches tapados y el profesor la vista
completa.}
\label{fig:archdinov2}
\end{figure}

\subsection{Lo común a los tres}

\begin{table}[H]
\centering
\small
\begin{tabularx}{\textwidth}{l>{\raggedright\arraybackslash}X}
\toprule
Ajuste & Valor \\
\midrule
Optimizador & AdamW, tasa base $10^{-4}$ \\
Tasa por época & subida lineal 20 épocas ($10^{-4}\cdot(e+1)/20$), luego coseno hasta $10^{-6}$ en la época 100 \\
Weight decay & VICReg $10^{-6}$; BYOL 0.01; DINOv2 de 0.04 a 0.2 (coseno por paso), sin aplicarlo a sesgos ni normalizaciones \\
Lote, épocas, pasos & 16 sujetos, 100 épocas, 2300 pasos \\
Precisión & bf16 mixta; pérdidas y cabeza DINO en precisión simple \\
Semillas & 1, 2 y 3 \\
Qué se evalúa & el estado final de la época 100, sin elegir época (sin \emph{early stopping}) \\
\bottomrule
\end{tabularx}
\end{table}

% ============================================================
\section{Qué se compara: supervisado, SSL, sin SSL y aleatorio}
\label{sec:conditions}
% ============================================================

\begin{tldrbox}
Las cuatro condiciones usan \emph{el mismo encoder y la misma entrada}. Cambia qué se entrena y con
qué. Solo el supervisado ajusta el encoder con las etiquetas; las otras tres lo congelan y entrenan
encima la misma regresión logística.
\end{tldrbox}

\begin{figure}[H]
\centering
\begin{tikzpicture}[x=1cm,y=1cm]
\node[lane=supc] at (0,0) {\textbf{Supervisado completo}\\[2pt]
  Pesos ImageNet inflados $\rightarrow$ encoder + lineal $\rightarrow$ 1 logit.\\[2pt]
  Se entrena \textbf{todo} con BCE en los 266 sujetos AD/HC de train.\\[2pt]
  Época elegida por AUROC de val.\\[2pt]
  Puntaje = logit calibrado.};
\node[lane=steelblue] at (4.1,0) {\textbf{SSL + probe}\\[2pt]
  Pesos ImageNet inflados $\rightarrow$ 100 épocas de SSL \textbf{sin etiquetas} (361 sujetos).\\[2pt]
  Encoder \textbf{congelado}.\\[2pt]
  Probe lineal con las etiquetas.\\[2pt]
  Puntaje = probabilidad del probe.};
\node[lane=softgray] at (8.2,0) {\textbf{Sin SSL + probe}\\[2pt]
  Pesos ImageNet inflados, \textbf{sin ningún entrenamiento}.\\[2pt]
  Encoder congelado.\\[2pt]
  El mismo probe.\\[2pt]
  Mide cuánto aporta el SSL sobre su punto de partida.};
\node[lane=softgray] at (12.3,0) {\textbf{Aleatorio + probe} (ref.)\\[2pt]
  Pesos \textbf{aleatorios}, sin entrenar.\\[2pt]
  Encoder congelado.\\[2pt]
  El mismo probe.\\[2pt]
  Mide cuánto aporta todo (ImageNet + SSL) sobre nada.};
\end{tikzpicture}
\caption{Las cuatro condiciones. Las tres primeras son las barras de la Figura~\ref{fig:metrics}, en
ese orden; la aleatoria aparece solo en las tablas.}
\label{fig:conditions}
\end{figure}

\begin{warnbox}[``Sin SSL'' no es ``aleatorio'']
En las figuras, \textbf{sin SSL} es el encoder con los pesos de ImageNet inflados, el mismo punto del
que arranca el SSL. Es la comparación justa para preguntar ``¿qué agregó el SSL?''. El encoder con
\textbf{pesos aleatorios} es otra referencia, más débil en principio, y responde ``¿qué agregó todo
el preentrenamiento?''. Los pesos sin SSL no dependen de la semilla (solo el \emph{mask token} del
ViT, que no se usa al evaluar), así que una sola corrida vale para las tres semillas.
\end{warnbox}

\subsection{El supervisado, en detalle}

\texttt{scripts/3\_train/train\_supervised.py}. Mismo encoder, mismos pesos iniciales, misma entrada y la
primera vista de la misma receta de augmentation.\footnote{Excepción: el supervisado de DINOv2 usó
la deformación elástica calculada en GPU (\texttt{elastic\_gpu}, desplazamiento 1.50 vóxeles,
calibrada a la misma RMSE), no la de TorchIO del SSL.}
\begin{itemize}
\item \textbf{Cabeza}: una lineal del vector (768/512/384) a un solo logit.
\item \textbf{Pérdida}: entropía cruzada binaria con peso para AD igual a $\#\text{HC}/\#\text{AD}$
  (135/131) en los sujetos de train.
\item \textbf{Optimización}: AdamW con tasa $10^{-4}$ en el encoder y $10^{-3}$ en la cabeza, weight
  decay 0.01, 5 épocas de subida y coseno por paso hasta la época 100; lotes de 16 (se descarta el
  lote incompleto: 16 pasos por época); una sesión al azar por sujeto en cada época.
\item \textbf{Selección}: tras cada época se mide el AUROC en los 63 de val (puntaje del sujeto =
  logit promedio de sus sesiones) y se guarda la mejor época; se corta tras 20 épocas sin mejora.
  Épocas elegidas: ConvNeXt 49, 12, 15; ResNet 13, 7, 6; ViT 26, 18, 51.
\item \textbf{Calibración}: una regresión logística balanceada sobre los logits de sus propios
  sujetos de entrenamiento (Platt). No cambia el AUROC de test; corrige el umbral de 0.5 del F1 y
  hace comparables los pliegues de la CV.
\item \textbf{CV}: 5 pliegues estratificados sobre los mismos 329 sujetos de train+val que usa el
  probe; cada pliegue entrena hasta la época elegida, se calibra con sus sujetos y predice los suyos.
\end{itemize}

% ============================================================
\section{Cómo se calcula cada número}
\label{sec:metrics}
% ============================================================

\begin{tldrbox}
Cada sujeto se resume en un vector (cubo limpio, encoder congelado, promedio de sus sesiones). Una
regresión logística sobre esos vectores da una probabilidad de AD por sujeto, y con ella se calculan
AUROC, F1 y AUPR, en validación cruzada (329 sujetos) y en test (84).
\end{tldrbox}

\subsection{Del volumen al vector del sujeto}

\begin{enumerate}
\item Para cada sesión de un sujeto AD o HC se toma el \textbf{cubo central de $32^3$ sin
  augmentation} (con todos sus canales).
\item El encoder, en modo evaluación (BatchNorm con sus estadísticas acumuladas, sin dropout ni drop
  path, en precisión simple), da un vector: 768 (ConvNeXt), 512 (ResNet) o 384 (CLS del ViT).
\item Se \textbf{promedian los vectores de todas las sesiones del sujeto}. Los sujetos de train
  pueden tener varias; los de val y test, una.
\item Quedan 413 vectores: 266 de train, 63 de val y 84 de test. Etiqueta: AD $=1$, HC $=0$.
\end{enumerate}

\subsection{El probe}

Una tubería de dos pasos (scikit-learn): \textbf{estandarizar} cada dimensión (media 0, desviación 1,
con las medias y desviaciones de los sujetos de ajuste) y una \textbf{regresión logística} con
regularización L2 y clases balanceadas. Su salida es una probabilidad de AD entre 0 y 1. Es lineal:
solo puede trazar un hiperplano en el espacio de los vectores, así que mide qué tan
\emph{linealmente} separable es el diagnóstico en la representación.

\subsection{\texorpdfstring{CV $5\times5$}{CV 5x5} (329 sujetos de train + val)}

\begin{enumerate}
\item Se repite 5 veces ($r=0,\ldots,4$) una partición estratificada en 5 pliegues (semilla $r$).
\item En cada pliegue, el probe se ajusta con 4/5 de los sujetos (unos 263); su regularización $C$
  se elige entre $\{0.01, 0.1, 1, 10\}$ con una CV interna de 3 pliegues que maximiza el AUROC. Luego
  predice el 1/5 restante (unos 66), que no vio.
\item Al terminar los 5 pliegues, cada sujeto tiene una predicción \emph{fuera de muestra}. Con las
  329 se calcula AUROC, F1 y AUPR.
\item Se promedia sobre las 5 repeticiones y luego sobre las 3 semillas.
\end{enumerate}

\begin{warnbox}[CV y test no miden lo mismo]
266 de los 329 sujetos de la CV estuvieron en el entrenamiento SSL (sin etiquetas). El encoder los
``conoce'', aunque nunca vio su diagnóstico. El test son sujetos que el encoder nunca vio, pero con
84 sujetos sus intervalos son anchos ($\pm0.05$ a $\pm0.07$ de AUROC).
\end{warnbox}

\subsection{Test (84 sujetos)}

El probe se ajusta en los 266 sujetos de train con cada $C\in\{0.01,0.1,1,10\}$; se queda el $C$ con
mejor AUROC en los 63 de val, y ese modelo predice una sola vez a los 84 de test.

\subsection{Las métricas}

\begin{itemize}
\item \textbf{AUROC}: la probabilidad de que un AD elegido al azar reciba un puntaje mayor que un HC
  elegido al azar. 0.5 es azar, 1 es separación perfecta. No depende de ningún umbral.
\item \textbf{F1-macro}: se clasifica como AD si la probabilidad es $\geq0.5$; se calcula el F1 de la
  clase AD y el de la clase HC y se promedian.
\item \textbf{AUPR} (\emph{average precision}): el área bajo la curva precisión--sensibilidad con AD
  como positivo. El azar vale la prevalencia de AD (0.49 en CV y en test).
\end{itemize}

\subsection{\texorpdfstring{Semillas, $\pm$ e intervalos de confianza}{Semillas, ± e intervalos de confianza}}

\begin{itemize}
\item El número de cada barra es la \textbf{media de las 3 semillas}; ``$\pm$'' es la desviación
  estándar entre semillas.
\item Las \textbf{diferencias} ($\Delta$) entre dos condiciones llevan un IC del 95\,\% por
  \emph{bootstrap pareado estratificado}: 2000 veces se remuestrean con reposición los sujetos AD y
  los HC por separado, se recalcula el AUROC de ambas condiciones \emph{sobre el mismo remuestreo}
  (promediando semillas) y se resta. El IC son los percentiles 2.5 y 97.5. En CV se usan las
  predicciones fuera de muestra promediadas en las 5 repeticiones.
\item \textbf{Significativo} ($^{*}$) quiere decir que el IC no toca el cero.
\end{itemize}

\begin{warnbox}[El CV del supervisado no es idéntico]
El probe repite 5 veces una CV de 5 pliegues; el supervisado hace una sola (5 entrenamientos por
semilla). El test (mismos 84 sujetos) es la comparación más limpia entre supervisado y SSL.
\end{warnbox}

\subsection{Monitor de colapso}
\label{sec:collapsedef}

Cada 2 épocas, sobre los 147 sujetos no vistos (63 de val y 84 de test, una sesión cada uno), el
encoder procesa el cubo limpio y dos vistas fijas de la receta del modelo:
\begin{itemize}
\item \textbf{Var sujeto / Var aug}: cuánto cambia el vector entre sujetos (cubos limpios) dividido
  por cuánto cambia entre las dos vistas del mismo sujeto. Sobre 1, el sujeto pesa más que la
  augmentation. Un modelo colapsado da lo mismo para todos los sujetos y el cociente cae a casi 0.
\item \textbf{Rango efectivo}: $\exp$ de la entropía del espectro de varianzas de los vectores. Vale
  1 si toda la variación cabe en una sola dirección.
\item \textbf{Top-1 de vistas}: fracción de primeras vistas cuya segunda vista más parecida (coseno)
  es la del mismo sujeto. Azar: $1/147$, un 0.7\,\%.
\end{itemize}

\subsection{\texorpdfstring{El espacio latente y qué es la $r$}{El espacio latente y qué es la r}}
\label{sec:rdef}

El análisis usa los 413 vectores de sujeto AD/HC de la semilla 1 (estado final), estandarizados.
\begin{itemize}
\item \textbf{PCA}: las 5 direcciones de mayor varianza. Cada sujeto tiene un puntaje en cada una.
\item \textbf{t-SNE} (perplejidad 30, inicio PCA, semilla 0): un mapa 2D que conserva vecindarios.
  Es solo visual: las distancias grandes y el tamaño de los grupos no significan nada.
\item \textbf{Puntaje del probe fuera de muestra}: regresión logística ($C$ entre 0.001 y 10, CV
  interna de 5) con 5 pliegues sobre los 413 sujetos. Por eso su AUROC (0.858, 0.858, 0.832) difiere
  algo del de la Sección~\ref{sec:results}.
\item \textbf{$\rho$ de cada componente}: correlación de Spearman (por rangos) entre el puntaje de
  una componente principal y una variable. El \textbf{signo de una componente es arbitrario}: PCA
  puede devolver el mismo eje apuntando hacia un lado o hacia el otro. Por eso se lee $|\rho|$.
\item \textbf{La $r$ (``embedding completo'')}: para cada variable, una regresión Ridge (con su
  regularización elegida por CV interna) aprende a predecirla desde el vector completo, fuera de
  muestra (5 pliegues). $r$ es la correlación de Spearman entre esa predicción y el valor real.
  Mide \emph{cuánto de esa variable se puede leer linealmente del embedding}, en sujetos que el
  modelo de lectura no vio.
\item \textbf{Para el diagnóstico}, $r$ es la Spearman entre el puntaje fuera de muestra del probe y
  la etiqueta 0/1. Con una etiqueta binaria y una prevalencia $p$, esa correlación es una función
  del AUROC:
  \[ r \approx \sqrt{12\,p(1-p)}\,(\mathrm{AUROC}-0.5) \approx 1.73\,(\mathrm{AUROC}-0.5)
     \quad (p=0.49). \]
  Así, $r=0.62$ corresponde a AUROC 0.86 (verificado en los datos), y con separación perfecta $r$ no
  pasa de $\approx0.87$. \textbf{Una $r$ de diagnóstico no se compara directamente con la $r$ de una
  variable continua}, que sí puede llegar a 1.
\end{itemize}

\textbf{Variables}: edad y años de educación (CSV clínico); volumen del hipocampo (etiquetas 17 y
53), de los ventrículos (4, 43, 5 y 44) y del cerebro (vóxeles $>0.1$), en espacio MNI y en cm$^3$;
intensidad media del cubo del hipocampo; y ruido de imagen: la dispersión robusta
($1.4826\times$MAD) del residuo de alta frecuencia (volumen menos su suavizado gaussiano) dentro de la
sustancia blanca erosionada de la resonancia completa, un tejido homogéneo.

% ============================================================
\section{Resultados esperados}
\label{sec:expected}
% ============================================================

\begin{tldrbox}
La expectativa inicial va en la dirección correcta, con cuatro precisiones: frente al supervisado lo
esperable es \emph{empatar}, no superar; ``sin SSL'' y ``aleatorio'' son dos referencias distintas;
la separación en el espacio latente es un gradiente y se mide con números, no a ojo; y las curvas de
pérdida SSL no se leen como las de un clasificador.
\end{tldrbox}

\begin{table}[H]
\centering
\small
\begin{tabularx}{\textwidth}{>{\raggedright\arraybackslash}p{3.6cm}>{\raggedright\arraybackslash}X>{\raggedright\arraybackslash}p{3.9cm}}
\toprule
Expectativa inicial & Versión precisa (lo que conviene esperar) & Cómo se mide \\
\midrule
Alcanzar o superar al supervisado completo &
  \textbf{No inferioridad}: el SSL con probe lineal empata con el supervisado (diferencia dentro del
  ruido) y nunca queda significativamente debajo. Superarlo es plausible cuando hay \emph{pocas}
  etiquetas; con 266 sujetos etiquetados y un supervisado que usa la misma augmentation y elige su
  mejor época en val, el empate es lo razonable. &
  $\Delta$ AUROC (SSL $-$ supervisado) con IC 95\,\%, sobre todo en test. \\
\addlinespace
Superar definitivamente al baseline sin SSL &
  Son \textbf{dos} expectativas. (a) Superar al mismo encoder \emph{sin SSL} (pesos ImageNet): el SSL
  adapta la red a resonancias, así que debería ganar. (b) Superar al encoder \emph{aleatorio}:
  debería ganar con más margen. &
  $\Delta$ AUROC con IC frente a cada referencia, en CV y en test. \\
\addlinespace
Separación lineal visible en los embeddings &
  Un \textbf{gradiente} AD--HC con solapamiento, no dos grupos: el Alzheimer es un continuo y un cubo
  de 64\,mm ve solo parte de la atrofia. Lo visual (PCA, t-SNE) es cualitativo; lo cuantitativo es el
  AUROC del probe fuera de muestra (bastante sobre 0.5) y que alguna componente principal ordene el
  diagnóstico. &
  AUROC fuera de muestra; $|\rho|$ de la 1.ª componente con AD. \\
\addlinespace
Correlación positiva de la división del espacio con el diagnóstico &
  El \textbf{signo} de una componente es arbitrario: se espera $|\rho|$ alto, no $\rho>0$. El puntaje
  del probe sí es positivo por construcción (AD $=1$). Además conviene esperar que el eje AD también
  correlacione con la edad (el confusor del split). &
  $|\rho|$ por componente; $r$ del diagnóstico y de la edad. \\
\addlinespace
Curvas: train y val bajan juntas, val algo arriba, sin sobreajuste ni inestabilidad &
  \textbf{Para el supervisado}, exactamente eso. \textbf{Para el SSL}, la pérdida debe bajar y
  estabilizarse, sin divergir y sin colapso; pero la de validación usa otros pares (el cubo y su
  espejo en BYOL y DINOv2) y otro modo de la red (BatchNorm en evaluación), así que puede quedar
  arriba, abajo o subir sin que el encoder empeore. El sobreajuste que importa se mide en sujetos no
  vistos: el probe en test y el monitor de colapso. &
  Curvas de pérdida; Var sujeto / Var aug, rango efectivo y top-1; métricas de 30 frente a 100
  épocas. \\
\bottomrule
\end{tabularx}
\caption{Expectativas, su versión precisa y cómo se mide cada una.}
\label{tab:expected}
\end{table}

\subsection{Lo esperado frente a lo obtenido}

\begin{table}[H]
\centering
\small
\begin{tabularx}{\textwidth}{>{\raggedright\arraybackslash}X ccc}
\toprule
Expectativa (versión precisa) & \textcolor{vicregc}{\textbf{VICReg}} & \textcolor{byolc}{\textbf{BYOL}} & \textcolor{dinoc}{\textbf{DINOv2}} \\
\midrule
Empata con el supervisado en test & \si{} ($-0.015$) & \si{} ($-0.012$) & \si{} ($+0.009$) \\
Empata o gana al supervisado en CV & \si{} ($+0.016$) & \si{} ($+0.026$\sig) & \si{} ($+0.002$) \\
Supera al encoder sin SSL en CV & \si{} ($+0.028$\sig) & \no{} ($-0.015$) & \si{} ($+0.047$\sig) \\
Supera al encoder sin SSL en test & \parcial{} ($+0.028$, n.s.) & \parcial{} ($+0.030$, n.s.) & \si{} ($+0.106$\sig) \\
Supera al encoder aleatorio (CV / test) & \si{} / \si{} & \no{} / \no{} (empata) & \si{} / \parcial{} \\
AUROC fuera de muestra en el espacio latente $\geq0.80$ & \si{} (0.858) & \si{} (0.858) & \si{} (0.832) \\
1.ª componente ordena el diagnóstico ($|\rho|\geq0.4$) & \si{} (0.49) & \si{} (0.59) & \si{} (0.56) \\
Sin colapso (Var sujeto/Var aug $\gg1$) & \si{} (47) & \si{} (57) & \si{} (17) \\
Pérdida SSL de val. sigue a la de train & \si{} & \no{} (sube desde la ép.\ 15) & \si{} (brecha desde la ép.\ 20) \\
Supervisado sin sobreajuste & \si{} & \no{} (época 6--13) & \parcial{} \\
\bottomrule
\end{tabularx}
\caption{$\Delta$ AUROC con IC 95\,\% por bootstrap pareado; \sig{} = significativo; n.s. = el IC
incluye el cero. ``En parte'' = el signo va en la dirección esperada sin significancia. Detalle en la
Sección~\ref{sec:results}.}
\label{tab:scorecard}
\end{table}

\begin{ideabox}[Tres sorpresas que vale la pena contar]
\begin{itemize}[leftmargin=12pt]
\item \textbf{Una ResNet-18 3D aleatoria ya separa bastante}: AUROC 0.840 en CV y 0.835 en test con
  pesos al azar y solo el probe. La atrofia del hipocampo (más líquido oscuro, menos tejido) es una
  señal tan fuerte que hasta filtros aleatorios la recogen. Por eso BYOL no supera ni al aleatorio
  ni a su inicio ImageNet: no hay mucho margen.
\item \textbf{El ViT con pesos de ImageNet arranca peor que con pesos aleatorios} (0.778 frente a 0.797
  en CV; 0.727 frente a 0.802 en test). El SSL lo lleva al nivel de las CNN. Ver la explicación
  probable en la Sección~\ref{sec:vitinflate}.
\item \textbf{100 épocas no mejoran a 30}: ninguna diferencia supera $\pm0.008$ de AUROC.
\end{itemize}
\end{ideabox}

% ============================================================
\section{Resultados}
\label{sec:results}
% ============================================================

\subsection{Métricas de los campeones}

\begin{figure}[H]
\centering
\includegraphics[width=\textwidth]{metricas_100_puntos.png}
\caption{AUROC, F1-macro y AUPR en CV (arriba) y test (abajo). Cuadrado: supervisado; círculo de
color: SSL; rombo vacío: sin SSL. El número junto a cada marcador es la media de las 3 semillas; la
barra horizontal es $\pm1$ desviación estándar y los puntos pequeños, las semillas. Cómo se calcula
cada métrica: Sección~\ref{sec:metrics}.}
\label{fig:metrics}
\end{figure}

\begin{table}[H]
\centering
\scriptsize
\setlength{\tabcolsep}{3.2pt}
\begin{tabular}{llcccccc}
\toprule
 & & \multicolumn{3}{c}{CV $5\times5$ (329 sujetos)} & \multicolumn{3}{c}{Test (84 sujetos)} \\
\cmidrule(lr){3-5}\cmidrule(lr){6-8}
Modelo & Condición & AUROC & F1-macro & AUPR & AUROC & F1-macro & AUPR \\
\midrule
\multirow{4}{*}{\textcolor{vicregc}{\textbf{VICReg}}}
 & Supervisado & 0.864\,$\pm$\,0.014 & 0.793\,$\pm$\,0.003 & 0.856\,$\pm$\,0.013 & \textbf{0.849}\,$\pm$\,0.008 & \textbf{0.764}\,$\pm$\,0.018 & \textbf{0.831}\,$\pm$\,0.009 \\
 & SSL 100 & \textbf{0.877}\,$\pm$\,0.003 & \textbf{0.796}\,$\pm$\,0.003 & \textbf{0.867}\,$\pm$\,0.011 & 0.834\,$\pm$\,0.007 & 0.756\,$\pm$\,0.008 & 0.798\,$\pm$\,0.021 \\
 & Sin SSL & 0.848 & 0.771 & 0.838 & 0.806 & 0.738 & 0.795 \\
 & \textit{Aleatorio} & \textit{0.808\,$\pm$\,0.011} & \textit{0.742\,$\pm$\,0.013} & \textit{0.809\,$\pm$\,0.015} & \textit{0.744\,$\pm$\,0.025} & \textit{0.677\,$\pm$\,0.031} & \textit{0.723\,$\pm$\,0.026} \\
\midrule
\multirow{4}{*}{\textcolor{byolc}{\textbf{BYOL}}}
 & Supervisado & 0.827\,$\pm$\,0.004 & 0.742\,$\pm$\,0.015 & 0.813\,$\pm$\,0.006 & \textbf{0.850}\,$\pm$\,0.016 & 0.765\,$\pm$\,0.025 & \textbf{0.812}\,$\pm$\,0.024 \\
 & SSL 100 & 0.846\,$\pm$\,0.018 & 0.765\,$\pm$\,0.003 & 0.819\,$\pm$\,0.039 & 0.838\,$\pm$\,0.020 & \textbf{0.785}\,$\pm$\,0.021 & 0.799\,$\pm$\,0.043 \\
 & Sin SSL & \textbf{0.860} & \textbf{0.778} & \textbf{0.832} & 0.808 & 0.738 & 0.784 \\
 & \textit{Aleatorio} & \textit{0.840\,$\pm$\,0.008} & \textit{0.769\,$\pm$\,0.008} & \textit{0.811\,$\pm$\,0.014} & \textit{0.835\,$\pm$\,0.011} & \textit{0.764\,$\pm$\,0.037} & \textit{0.788\,$\pm$\,0.015} \\
\midrule
\multirow{4}{*}{\textcolor{dinoc}{\textbf{DINOv2}}}
 & Supervisado & \textbf{0.833}\,$\pm$\,0.019 & \textbf{0.760}\,$\pm$\,0.016 & 0.824\,$\pm$\,0.021 & 0.825\,$\pm$\,0.026 & \textbf{0.744}\,$\pm$\,0.018 & \textbf{0.804}\,$\pm$\,0.034 \\
 & SSL 100 & 0.832\,$\pm$\,0.005 & 0.748\,$\pm$\,0.002 & \textbf{0.825}\,$\pm$\,0.005 & \textbf{0.833}\,$\pm$\,0.013 & 0.740\,$\pm$\,0.009 & 0.794\,$\pm$\,0.014 \\
 & Sin SSL & 0.778 & 0.709 & 0.750 & 0.727 & 0.677 & 0.645 \\
 & \textit{Aleatorio} & \textit{0.797\,$\pm$\,0.014} & \textit{0.733\,$\pm$\,0.013} & \textit{0.764\,$\pm$\,0.018} & \textit{0.802\,$\pm$\,0.032} & \textit{0.734\,$\pm$\,0.007} & \textit{0.801\,$\pm$\,0.041} \\
\bottomrule
\end{tabular}
\caption{Todas las métricas: media $\pm$ desviación estándar de 3 semillas (sin SSL es una corrida).
En negrita, la mejor de las tres condiciones principales en cada columna. En cursiva, la referencia
con pesos aleatorios.}
\label{tab:allmetrics}
\end{table}

\subsection{Diferencias con intervalo de confianza}

\begin{figure}[H]
\centering
\includegraphics[width=0.82\textwidth]{delta_100.png}
\caption{$\Delta$ AUROC del SSL frente al mismo encoder sin SSL (círculo lleno) y frente al
supervisado (cuadrado vacío), con IC del 95\,\% por bootstrap pareado. Un intervalo que no toca el
cero es una diferencia significativa.}
\label{fig:delta}
\end{figure}

\begin{table}[H]
\centering
\footnotesize
\begin{tabular}{lcccc}
\toprule
$\Delta$ AUROC (CV / test) & SSL $-$ sin SSL & SSL $-$ supervisado & Supervisado $-$ sin SSL & SSL $-$ aleatorio \\
\midrule
\textcolor{vicregc}{\textbf{VICReg}} & $+.028^{*}$ / $+.028$ & $+.016$ / $-.015$ & $+.012$ / $+.043$ & $+.064^{*}$ / $+.090^{*}$ \\
\textcolor{byolc}{\textbf{BYOL}} & $-.015$ / $+.030$ & $+.026^{*}$ / $-.012$ & $-.040^{*}$ / $+.043$ & $+.010$ / $+.003$ \\
\textcolor{dinoc}{\textbf{DINOv2}} & $+.047^{*}$ / $+.106^{*}$ & $+.002$ / $+.009$ & $+.045^{*}$ / $+.097^{*}$ & $+.033^{*}$ / $+.031$ \\
\bottomrule
\end{tabular}
\caption{$^{*}$ = IC 95\,\% enteramente de un lado del cero (BYOL supervisado $-$ sin SSL en CV es
significativamente \emph{negativo}). Intervalos completos en
\nolinkurl{campaign_ssl3d_runpod/split361_100/REPORT.md}.}
\label{tab:deltas}
\end{table}

\begin{ideabox}[Cómo leerlo]
\begin{itemize}[leftmargin=12pt]
\item \textbf{Frente al supervisado}, el SSL empata en test en los tres (de $-0.015$ a $+0.009$) y
  empata o gana en CV. Con 266 sujetos etiquetados, congelar el encoder SSL y entrenar solo una
  regresión logística rinde lo mismo que entrenar la red entera con las etiquetas.
\item \textbf{Frente al encoder sin SSL}, DINOv2 gana claro (parte del peor inicio), VICReg gana en CV
  y BYOL no gana.
\item \textbf{Frente al aleatorio}, VICReg gana claro; DINOv2 gana en CV; BYOL empata.
\end{itemize}
\end{ideabox}

\subsection{El espacio latente}
\label{sec:latent}

Los embeddings salen de los modelos de 100 épocas (semilla 1, estado final), sin reentrenar nada.
Definiciones en la Sección~\ref{sec:rdef}.

\begin{figure}[H]
\centering
\includegraphics[width=0.92\textwidth]{latent_100.png}
\caption{Cada punto es un sujeto (413). Fila 1: las dos componentes principales. Fila 2: t-SNE.
Fila 3: el puntaje del probe para cada sujeto, calculado sin que el probe lo haya visto; la línea
punteada es el umbral.}
\label{fig:latent}
\end{figure}

\textbf{Cómo leer la Figura~\ref{fig:latent}.} En PCA y t-SNE AD y HC no forman dos grupos, sino un
gradiente: los AD se concentran hacia un lado, con mucha mezcla en el medio. Es lo esperado
(Tabla~\ref{tab:expected}). La fila 3 es la que cuantifica la separación lineal: los dos histogramas
se separan con solapamiento, con AUROC fuera de muestra de 0.858 (VICReg), 0.858 (BYOL) y 0.832
(DINOv2), frente a 0.842, 0.849 y 0.767 de los mismos encoders sin SSL.

\begin{figure}[H]
\centering
\includegraphics[width=\textwidth]{latent_corr_100.png}
\caption{Correlación de Spearman de cada variable con las 5 primeras componentes ($\rho$; el signo es
arbitrario) y con su predicción lineal desde el vector completo ($r$, última columna de cada bloque).}
\label{fig:latentcorr}
\end{figure}

\begin{itemize}
\item \textbf{La 1.ª componente es un eje de atrofia}: ordena el volumen del hipocampo
  ($|\rho|$ 0.67--0.84), el de los ventrículos (0.68--0.81) y la edad (0.58--0.61), y con ellos el
  diagnóstico (0.49--0.59).
\item \textbf{El vector completo codifica casi perfecto la anatomía}: $r$ de 0.88--0.90 para el
  hipocampo y 0.82--0.88 para los ventrículos. La educación no aparece ($|r|\leq0.06$): es el
  control negativo.
\item \textbf{También codifica intensidad y ruido} ($r$ 0.67--0.93): el encoder describe toda la
  imagen, no solo lo útil para el diagnóstico.
\item \textbf{La $r$ del diagnóstico} (0.57--0.62) es la del AUROC $\approx$0.83--0.86, con techo de
  0.87; no hay que compararla de igual a igual con la $r$ de la anatomía.
\end{itemize}

\begin{figure}[p]
\centering
\includegraphics[width=0.88\textwidth]{latent_tsne_vars_100.png}
\caption{El mismo t-SNE de cada modelo, coloreado por cada variable (color = percentil). Una
variable codificada se ve como un gradiente suave; una no codificada, como sal y pimienta
(educación). En cada panel, la $r$ de esa variable; con marco de color, las que el embedding predice
al menos tan bien como el diagnóstico.}
\label{fig:tsnevars}
\end{figure}

\begin{figure}[H]
\centering
\includegraphics[width=\textwidth]{noise_checks_100.png}
\caption{Tres pruebas de que el eje AD se apoya en la anatomía y no en el ruido. 1: AUROC de test
con ruido gaussiano añadido al cubo (el probe se ajustó con cubos limpios). 2: ¿el vecino más
cercano de una sesión es otra sesión del mismo sujeto? 3: correlación parcial de Spearman del puntaje
AD con cada variable, descontando las otras (marca negra: la del diagnóstico real).}
\label{fig:noise}
\end{figure}

\begin{table}[H]
\centering
\small
\begin{tabular}{lccc}
\toprule
 & \textcolor{vicregc}{\textbf{VICReg}} & \textcolor{byolc}{\textbf{BYOL}} & \textcolor{dinoc}{\textbf{DINOv2}} \\
\midrule
AUROC del probe fuera de muestra (sin SSL entre paréntesis) & 0.86 (0.84) & 0.86 (0.85) & 0.83 (0.77) \\
Varianza de la 1.ª componente & 0.22 & 0.22 & 0.12 \\
$|\rho|$ 1.ª componente -- AD & 0.49 & 0.59 & 0.56 \\
$|\rho|$ 1.ª componente -- volumen del hipocampo & 0.67 & 0.84 & 0.76 \\
AUROC de test con ruido $\sigma=0.3$ & 0.84 & 0.84 & 0.84 \\
Otra sesión del mismo sujeto como vecino más cercano & 0.23 & 0.65 & 0.74 \\
Puntaje AD -- volumen del hipocampo (parcial) & 0.76 & 0.79 & 0.77 \\
Puntaje AD -- ruido (parcial; diagnóstico real: 0.10) & 0.23 & 0.23 & 0.39 \\
\bottomrule
\end{tabular}
\caption{Resumen del espacio latente a 100 épocas (semilla 1).}
\label{tab:latent}
\end{table}

\subsection{Curvas de entrenamiento}

\begin{figure}[H]
\centering
\includegraphics[width=\textwidth]{split361_100_curvas_vicreg.png}
\includegraphics[width=0.8\textwidth]{split361_100_curvas_byol.png}
\includegraphics[width=\textwidth]{split361_100_curvas_dinov2.png}
\caption{Pérdidas SSL en 100 épocas (continua: train; discontinua: validación; gruesa: media de 3
semillas).}
\label{fig:curves}
\end{figure}

\begin{itemize}
\item \textbf{VICReg}: la pérdida total baja en train y en val, con val ligeramente arriba: el patrón
  esperado. La invarianza de val queda \emph{debajo} de la de train porque en val las dos vistas
  se normalizan con las mismas estadísticas de BatchNorm. No hay sobreajuste ni inestabilidad.
\item \textbf{BYOL}: train baja hasta 0.02, y val baja hasta las épocas 13--17 y luego sube y se
  estabiliza en 0.24--0.53. No es lo esperado para una curva de clasificador, pero el diagnóstico del
  informe anterior mostró que viene de que las cabezas se evalúan en otro modo, y que a 100 épocas
  aparece una brecha pequeña (0.03--0.07) entre sujetos de train y no vistos. El probe no lo nota.
\item \textbf{DINOv2}: las tres pérdidas bajan hasta el final; val queda sobre train y la brecha se
  abre desde la época 20, pero val nunca sube. La entropía del profesor se estabiliza en 0.65 de su
  máximo: sin colapso.
\end{itemize}

\begin{figure}[H]
\centering
\includegraphics[width=\textwidth]{supervisado_curvas_100.png}
\caption{Supervisado: BCE de entrenamiento (continua, con vistas aumentadas y peso por clase) y de
validación (discontinua, 63 sujetos con cubo limpio, logit sin calibrar), una línea por semilla.
Punto: la época elegida por AUROC de validación. La ResNet-18 sobreajusta en pocas épocas: su BCE de
train cae a 0.05 y la de validación salta a 1--4 (predicciones demasiado seguras), aunque su AUROC
siga alto; por eso el early stopping elige las épocas 6--13.}
\label{fig:supcurves}
\end{figure}

\begin{figure}[H]
\centering
\includegraphics[width=0.86\textwidth]{supervised_100.png}
\caption{Supervisado: la otra mitad de la historia, el AUROC de validación por época.}
\label{fig:supauc}
\end{figure}

\subsection{Colapso}

\begin{figure}[H]
\centering
\includegraphics[width=0.86\textwidth]{collapse_100.png}
\caption{Monitor de colapso en los 147 sujetos no vistos (definiciones en la
Sección~\ref{sec:collapsedef}). Sobre 1, el sujeto pesa más que la augmentation; punteada: la corrida
de 30 épocas.}
\label{fig:collapse}
\end{figure}

\begin{table}[H]
\centering
\small
\begin{tabular}{lccc}
\toprule
Al final (época 100) & Var sujeto / Var aug & Rango efectivo & Top-1 de vistas \\
\midrule
\textcolor{vicregc}{\textbf{VICReg}} & 47.1 & 9.5 & 0.95 \\
\textcolor{byolc}{\textbf{BYOL}} & 56.7 & 13.2 & 0.99 \\
\textcolor{dinoc}{\textbf{DINOv2}} & 17.4 & 25.6 & 1.00 \\
\bottomrule
\end{tabular}
\caption{Ninguno colapsa: los tres reconocen a casi todos los sujetos no vistos entre 147 (azar:
0.007).}
\end{table}

\begin{warnbox}[Reconocer al sujeto no es reconocer la enfermedad]
Los tres aprenden muy bien \emph{quién} es cada sujeto, pero el probe mejora poco sobre sus puntos de
partida. La identidad del sujeto no basta para separar AD de HC.
\end{warnbox}

% ============================================================
\section{En resumen}
% ============================================================

\begin{enumerate}[leftmargin=16pt]
\item \textbf{Entrada}: cubos 3D de $32^3$ vóxeles (64\,mm) en posiciones fijas del espacio MNI; el
  hemisferio derecho espejado; 2 canales (hipocampos) para las CNN y 6 (hipocampo, ventrículo y
  corteza posteromedial, de cada lado) para el ViT.
\item \textbf{Pesos}: ImageNet-1k en 2D (torchvision y DeiT), ninguno médico, inflados a 3D con la
  regla de I3D; en el ViT, además, redimensionado del parche (FlexiViT) y posiciones 3D.
\item \textbf{Arquitecturas}: ConvNeXt-T 3D (768), ResNet-18 3D sin max-pool (512) y ViT-S 3D con 64
  parches de $8^3$ (CLS de 384). Solo ese vector se evalúa.
\item \textbf{Resultados}: el SSL empata con el supervisado completo; frente al encoder sin SSL gana
  DINOv2 (mucho), VICReg (en CV) y BYOL no; los embeddings ordenan el diagnóstico como un gradiente
  apoyado en la atrofia del hipocampo, con la edad como confusor.
\end{enumerate}

\begin{warnbox}[Límites que conviene decir en voz alta]
\begin{itemize}[leftmargin=12pt]
\item Una sola cohorte (OASIS3) y un solo split; falta validar en otra cohorte o escáner.
\item La edad separa AD de HC casi tanto como el diagnóstico.
\item Test de 84 sujetos: intervalos de $\pm0.05$--$0.07$ de AUROC.
\item La CV incluye sujetos que el SSL vio (sin etiquetas).
\item Aún no hay una curva con pocas etiquetas (10, 25, 50\,\% de train), que es donde el SSL debería
  separarse del supervisado; el código para hacerla existe
  (\texttt{evaluate\_label\_fractions}).
\end{itemize}
\end{warnbox}

\appendix
\section{Reproducir este documento}

\begin{itemize}
\item Figuras nuevas (\texttt{ingesta\_3d}, \texttt{volumenes\_3d}, \texttt{architectures\_100},
  \texttt{metricas\_100\_puntos}, \texttt{supervisado\_curvas\_100}):
  \texttt{python scripts/5\_reports/build\_report\_metodologia.py}.
\item Figuras del split y de las vistas: \texttt{python scripts/5\_reports/build\_report\_split361.py};
  métricas, curvas y colapso a 100 épocas: \texttt{python scripts/5\_reports/report\_split361\_100.py}; espacio
  latente: \texttt{python scripts/5\_reports/analyze\_latent\_split361.py} con \texttt{-{}-runs split361\_100
  -{}-out latent\_split361\_100 -{}-suffix \_100}.
\item Entrenamiento: \texttt{run\_all.py} con \texttt{-{}-epochs 100 -{}-identity -{}-collapse-every 2}
  y \nolinkurl{vicreg:geo:32:campaign_augval}, \nolinkurl{byol:noise:32:gn_dp0.3_dq0_wd0.01},
  \nolinkurl{dinov2:elastic:multi3}; supervisado con \texttt{scripts/3\_train/train\_supervised.py}; encoders
  sin entrenar con \texttt{train.py -{}-init-only}.
\item Compilar desde \texttt{docs/}: dos veces \texttt{pdflatex informe\_metodologia\_split361.tex}.
\end{itemize}

\end{document}
